% ======================================================================
\section{Introduction}\label{sec:intro}
% ======================================================================

At high energy the wave function of a hadron or nucleus contains a large number of soft gluons, and in a sufficiently dense system their density saturates~\cite{Gribov:1983ivg,Mueller:1985wy}. A description of this regime from first principles was proposed by McLerran and Venugopalan~\cite{McLerran:1993ni,McLerran:1993ka}: the fast (valence) partons of a large nucleus are treated as static sources of colour charge, with a density $\rho^a(\v x)$ in the transverse plane; the soft gluons are the classical Yang--Mills field radiated by these sources; and observables are obtained by averaging over the sources with a Gaussian weight. This picture has developed into the effective theory known as the Color Glass Condensate (CGC)~\cite{Iancu:2003xm,Weigert:2005us,Gelis:2010nm,Kovchegov:2012mbw}. The dependence of observables on the rapidity interval over which the soft modes are integrated out is governed by the BFKL equation~\cite{Kuraev:1977fs,Balitsky:1978ic} in the dilute regime and by the non-linear Balitsky--Kovchegov (BK) and JIMWLK equations~\cite{Balitsky:1995ub,Kovchegov:1999yj,Kovchegov:1999ua,Jalilian-Marian:1997qno,Jalilian-Marian:1997jx,Jalilian-Marian:1997ubg,Kovner:2000pt,Weigert:2000gi,Iancu:2000hn,Iancu:2001ad,Ferreiro:2001qy} in the dense regime, which replace the Gaussian weight by an evolving distribution of the sources.

The simplest observable that is sensitive both to the projectile and to the target wave functions is the single inclusive gluon spectrum in the collision of a dilute projectile with a dense target. At leading order (LO), for a gluon produced at mid rapidity, i.e.\ far in rapidity from both the projectile and the target, the spectrum is given by the formula of Kovchegov and Tuchin~\cite{Kovchegov:2001sc,Kovchegov:1998bi}: the gluon is emitted from the valence charges of the projectile in the Weizs\"acker--Williams field, and the emission amplitude, together with the colour rotation of the emitting charge, scatters eikonally on the target. The cross section is a product of the projectile colour charge correlator and of a target operator built from two adjoint Wilson lines; in the limit of a dilute target it reduces to the $k_T$-factorised form with the Lipatov vertex~\cite{Gribov:1983ivg,Kharzeev:2003wz,Blaizot:2004wu}. This is a result of the classical theory: it is quadratic in the charges of the projectile, and after the average over a colour-neutral projectile it does not depend on whether these charges are numbers or operators. Its leading-logarithmic rapidity evolution was obtained by Kovner and Lublinsky~\cite{Kovner:2006wr} (see also~\cite{Kovner:2006ge}) and, within the Schwinger--Keldysh formulation, by Gelis, Lappi and Venugopalan~\cite{Gelis:2008rw}: the projectile is to be evolved from its own rapidity down to the rapidity of the produced gluon, and the target from the rapidity of the produced gluon to its own rapidity.

Beyond leading order the situation is different. The original picture of classical fields of static, commuting colour charges with a Gaussian weight was designed for the leading-order problem and solves it; it is not by itself a complete description of gluon production at the next order. Three elements have to be added to it. First, the colour charges of the valence partons are operators, $[\rho^a(\v x),\rho^b(\v y)]=if^{abc}\rho^c(\v x)\delta^{(2)}(\v x-\v y)$~\cite{Kovner:2005pe,Kovner:2005nq,Kovner:2005jc,Kovner:2005en}. At the order $g^4$ of the next-to-leading order (NLO) correction the amplitudes contain products of three and four charges, and a commutator of two charges is of the same order as the two-charge terms that carry the evolution and the NLO coefficient. The order of the charges therefore matters: the commutators leave contact terms, and the totally symmetrised product of three charges brings in the three-point function of the projectile with the symmetric structure constants $d^{abc}$, which is odd under charge conjugation and absent for a Gaussian weight. Second, the logarithms of an NLO calculation are of several kinds: those of the rapidity evolution of the projectile and of the target, the collinear logarithms of the fragmentation of the produced gluon, which are governed by the DGLAP equations~\cite{Gribov:1972ri,Altarelli:1977zs,Dokshitzer:1977sg}, and the ultraviolet logarithms that make the coupling run. In the leading-order framework they are supplied by separate considerations; in a complete description all of them, with their scales and their overlaps, have to emerge from one calculation. Third, the terms of relative order $\alpha_s$ that are not enhanced by any logarithm have not been available for this observable.

Many of the ingredients have become available over the last years. The evolution equations themselves are known at NLO: the NLO BK equation~\cite{Balitsky:2007feb} and the NLO JIMWLK Hamiltonian~\cite{Balitsky:2013fea,Kovner:2013ona,Kovner:2014lca,Grabovsky:2013mba}, together with the running-coupling corrections to the LO kernels~\cite{Balitsky:2006wa,Kovchegov:2006vj,Kovchegov:2006wf}. In particular, Lublinsky and Mulian~\cite{Lublinsky:2016meo} computed the soft gluon light-cone wave function of a dilute projectile to order $g^3$ and derived the NLO JIMWLK Hamiltonian from it. For particle production, the NLO corrections to single inclusive hadron production at forward rapidity were computed in the hybrid formalism~\cite{Dumitru:2005gt,Altinoluk:2011qy} by Chirilli, Xiao and Yuan~\cite{Chirilli:2011km,Chirilli:2012jd}, and the structure of the rapidity subtraction and of the kinematic constraints in this framework was analysed in~\cite{Stasto:2013cha,Altinoluk:2014eka,Ducloue:2016shw}. Running-coupling corrections to inclusive gluon production at mid rapidity were obtained by Horowitz and Kovchegov~\cite{Horowitz:2010yg}, and the relation between $\beta_0$-dependent terms of NLO calculations and the running of the coupling was clarified in~\cite{Kovner:2023vsy}. The general structure of the NLO correction to inclusive spectra in the CGC was discussed by Gelis and Laidet~\cite{Gelis:2012ri}. Corrections of higher order in the projectile colour charge were computed in the classical approximation in~\cite{Chirilli:2015tea,Li:2021zmf,Li:2021ntm}, and their contribution to single inclusive production, odd in the transverse momentum, was analysed in~\cite{McLerran:2016snu,Kovchegov:2018jun}. A complete NLO calculation of the single inclusive gluon spectrum at mid rapidity has, to our knowledge, not been presented.

The purpose of this paper is to supply the missing pieces for this observable, in pure Yang--Mills theory. We work in light-cone perturbation theory (LCPT)~\cite{Lepage:1980fj,Brodsky:1997de}. The soft gluons of the projectile are generated from the valence charge density by a unitary evolution operator $\Omega$~\cite{Kovner:2005en}, which was constructed to order $g^2$ in~\cite{Radhakrishnan:2026soft}. It is written as a normal-ordered series in soft gluon creation and annihilation operators, with coefficients that are functionals of the non-commuting charge density $\rho$ and that are fixed by matching to LCPT and by unitarity. The NLO cross section requires in addition the coefficients that create or absorb one gluon at order $g^3$. The one that acts on the vacuum is the one-gluon component of the wave function of~\cite{Lublinsky:2016meo}, which we quote; the order of the charges in its three-charge term, which is essential here, is determined in this paper, and the other coefficient follows from unitarity. The observable is the expectation value of the number operator of dressed gluons in the scattered projectile state~\cite{Kovner:2006wr}. The calculation therefore reduces to an algebraic evaluation of products of the coefficients of $\Omega$ and of Wilson lines, in which the order of the charges is kept, followed by longitudinal and transverse integrations. No average over the projectile is taken until the end, and no model for it is assumed.

Our main results are the following.
\begin{itemize}
\item The cross section at order $g^4$ decomposes into three groups of terms: the interference of the LO amplitude with the one-gluon amplitude of order $g^3$ (Group~I); the squares of the parts of the two-gluon amplitude of order $g^2$, with one of the gluons integrated over (Group~II); and the interferences of different parts of the two-gluon amplitude (Group~III). We list all fifteen products of coefficients, the \emph{rows}, with their longitudinal integration ranges and their order in $\rho$ (Section~\ref{sec:cross}).
\item All terms enhanced by $\log(\vee/\Lambda)$, where $\Lambda<p^+<\vee$ is the range of the longitudinal momenta of the soft modes, are the leading-logarithmic evolution of the LO result. The logarithm splits at the longitudinal momentum $k^+$ of the measured gluon into $\log(\vee/k^+)+\log(k^+/\Lambda)$; the coefficient of $\log(\vee/k^+)$ is exactly the BFKL evolution of the projectile charge correlator inserted into the LO formula, and the coefficient of $\log(k^+/\Lambda)$ is exactly the JIMWLK Hamiltonian acting on the target operator of the LO formula. The contributions quartic in $\rho$, all of which carry $\log(\vee/\Lambda)$, cancel identically for commuting charges, and the totally symmetrised cubic terms that carry a logarithm cancel as well; with non-commuting charges, their reordering produces part of the two-charge terms of the evolution (Section~\ref{sec:rapidity}).
\item The remaining terms are a collinear term, associated with the fragmentation of the measured gluon and governed by the gluon--gluon splitting function, a running-coupling term, and finite terms. We give an account of all poles of dimensional regularisation: the ultraviolet poles of the virtual corrections to the measured gluon cancel among themselves, the coupling is renormalised by the normalisation of the dressed measured gluon, and the collinear poles combine into the full $\MS$ splitting function and are absorbed into the fragmentation function. The scheme is conventional dimensional regularisation with the $\MS$ coupling and the $\MS$ number-density fragmentation function (Section~\ref{sec:dglap_rc}).
\item We give the finite terms row by row, after the cancellations that are exact, for an arbitrary projectile, with every product of three charges written as its totally symmetrised product plus two-charge contact terms, and averaged over the projectile, in colour-singlet form. All colour structures are written, as at leading order, through differences of Wilson lines, and all longitudinal integrals with plus prescriptions in a momentum fraction. After the average the terms cubic in $\rho$ do not vanish: their symmetrised part, odd in $\v k$, is proportional to the symmetrised ($d^{abc}$) three-point function of the projectile, which is kept general, and to C-odd structures of the target, while the commutators of the charges leave contact terms, even in $\v k$, which are present for any projectile. The large-distance tails of individual terms cancel in pairs of a virtual and a real term, and we state which terms have to be integrated together (Section~\ref{sec:finite}, Appendices~\ref{app:finite} and~\ref{app:tails}).
\item In the limit in which also the target is dilute the real emission reproduces the two-gluon production cross section in quasi-multi-Regge kinematics~\cite{Fadin:1998py,Bartels:2006hg}, and for a point-like projectile the virtual correction has the gluon Regge trajectories as coefficients of the two rapidity logarithms, the one on the target side together with a soft--collinear term. The comparison of the virtual correction with the one-loop central-emission vertex, which requires the real emission with matched regulators, has not been completed (Section~\ref{sec:checks}).
\end{itemize}
The result can be written schematically as
\begin{align}\label{eq:schematic}
\frac{d^3N}{d^3k} &= \frac{d^3N_{\rm LO}}{d^3k}
+ \log\frac{\vee}{k^+}\,\frac{d^3N_{\rm BFKL}}{d^3k}
+ \log\frac{k^+}{\Lambda}\,\frac{d^3N_{\rm JIMWLK}}{d^3k}\notag\\
&\quad+ \frac{d^3N_{\rm DGLAP}}{d^3k}
+ \frac{d^3N_{\rm RC}}{d^3k}
+ \frac{d^3N_{\rm finite}}{d^3k}\,,
\end{align}
where every term on the right-hand side except the first is of relative order $\alpha_s$. The evolution of the projectile, the evolution of the target, the collinear evolution of the final state, the running of the coupling and the finite terms, with the charges of the projectile treated as operators throughout, come out of a single calculation. In this sense the present work completes the NLO description of this observable within the eikonal, dilute-projectile approximation, up to the open points listed in Section~\ref{sec:openpoints}. The qualification is essential. The finite terms are given in coordinate space, and no numerical evaluation of them is given; the comparison of the finite part of the virtual correction with the one-loop central-emission vertex in the dilute limit has not been completed; vanishing energy denominators, which occur in one coefficient of $\Omega$, are defined as principal values, and terms on the energy shell by which other prescriptions differ are not included. The paper does not address the collision of two dense objects, corrections beyond the eikonal approximation, or quarks.

The paper is organised as follows. In Section~\ref{sec:setup} we fix the conventions and review the light-cone Hamiltonian in the eikonal approximation, light-cone perturbation theory with operator-valued matrix elements, and the definition of $\Omega$. In Section~\ref{sec:omega} the coefficients of $\Omega$ are determined: those of the dressed vacuum are quoted from~\cite{Lublinsky:2016meo}, those of the dressed one-gluon state are computed, and the remaining ones follow from unitarity. In Section~\ref{sec:cross} we define the observable, derive the LO spectrum, and organise the cross section at order $g^4$ into Groups and rows. Section~\ref{sec:rapidity} is devoted to the rapidity logarithms and their identification with the BFKL and JIMWLK evolution. The poles, the collinear and the running-coupling terms are the subject of Section~\ref{sec:dglap_rc}, and the finite terms of Section~\ref{sec:finite}; a summary of the complete result is given in Section~\ref{sec:result_summary}. Section~\ref{sec:checks} lists the checks, with what each of them does and does not test, and the limits. Section~\ref{sec:discussion} contains the discussion and the list of open points, and Section~\ref{sec:conclusions} our conclusions. The notation is summarised in Table~\ref{tab:notation}. Appendices~\ref{app:hamiltonian}--\ref{app:rows} contain the details of the construction of $\Omega$ and of the cross section in terms of its coefficients; Appendix~\ref{app:finite} the results for an arbitrary projectile; Appendix~\ref{app:tails} the analysis of the large-distance behaviour and of the integrability of the finite terms; Appendix~\ref{app:distributions} the prescriptions; and Appendix~\ref{app:fourier} the Fourier transforms.

\begingroup
\footnotesize
\renewcommand{\arraystretch}{1.2}
\begin{longtable}{@{}>{\raggedright\arraybackslash}p{0.21\textwidth}>{\raggedright\arraybackslash}p{0.55\textwidth}>{\raggedright\arraybackslash}p{0.17\textwidth}@{}}
\caption{Notation. The last column gives the place of the definition.}\label{tab:notation}\\
\toprule
symbol & meaning & defined in\\
\midrule
\endfirsthead
\toprule
symbol & meaning & defined in\\
\midrule
\endhead
\bottomrule
\endfoot
\multicolumn{3}{@{}l}{\emph{Kinematics and cut-offs}}\\
$x^\pm$, $\v x$, $\int_x$ & light-cone coordinates $(x^0\pm x^3)/\sqrt2$; transverse vector; $\int d^2x$ & Sec.~\ref{sec:conventions}\\
$k=(k^+,\v k)$ & momentum of the measured gluon & Eq.~\eqref{eq:observable}\\
$p^+$ & longitudinal momentum of the second gluon (integrated over, or virtual) & Table~\ref{tab:rows}\\
$\vee$, $\Lambda$ & upper and lower cut-off of the soft modes, $\Lambda<p^+<\vee$; $\vee=\Lambda e^{\delta Y}$ & Eq.~\eqref{eq:softwindow}\\
$L$, $\ell_\vee$ & $\log(k^+/\Lambda)$; $\log(1-k^+/\vee)$, a power correction & Eq.~\eqref{eq:dg_B}\\
$\xi$ & momentum fraction: $p^+/k^+$ in all rows except Row~II.IV; $k^+/(p^++k^+)$ in Row~II.IV and in the DGLAP term & Eqs.~\eqref{eq:xi_def}, \eqref{eq:xi_defIV}\\
$\epsilon$, $\mu$ & $d=2-\epsilon$ transverse dimensions; $\MS$ scale ($\mu_R$, $\mu_F$ when distinguished) & Sec.~\ref{sec:dg_poles}\\
$c_0$ & $2e^{-\gamma_E}$ & Eq.~\eqref{eq:dg_B}\\
$\PV$ & principal value & Eq.~\eqref{eq:PV}\\[2pt]
\multicolumn{3}{@{}l}{\emph{Projectile}}\\
$\rho^a(\v x)$ & valence colour charge density, an operator & Eqs.~\eqref{eq:rhodef}, \eqref{eq:rhoalgebra}\\
$\langle\dots\rangle$ & average over the valence state of the projectile, taken last & Eq.~\eqref{eq:observable}\\
$n(\v x,\v y)$ & two-point function $\langle\rho^f(\v x)\rho^f(\v y)\rangle$ & Eqs.~\eqref{eq:singlet}, \eqref{eq:A1n}\\
$\sigma(\v x,\v y,\v z)$ & $\langle d^{efg}\{\rho^e(\v x)\rho^f(\v y)\rho^g(\v z)\}_{\rm sym}\rangle$, symmetrised three-point function, C-odd & Eq.~\eqref{eq:A2d}\\
$\{\dots\}_{\rm sym}$ & average of a product of charges over all orderings & Eq.~\eqref{eq:symord}\\[2pt]
\multicolumn{3}{@{}l}{\emph{Evolution operator and its coefficients}}\\
$a^a_i$, $a^{\dagger a}_i$ & soft gluon operators, colour $a$, transverse polarisation $i$ & Eq.~\eqref{eq:aadagger}\\
$\Omega$ & unitary evolution operator of the soft modes & Eq.~\eqref{eq:OmegaExp}\\
$\mf N=1+\mf N_2$ & coefficient of the identity; $\mf N_2=O(g^2\rho^2)$ & Eqs.~\eqref{eq:N2mom}, \eqref{eq:N2X}\\
$\mf A\equiv\mf A^{(1)}_2$ & emission of one gluon at order $g$ (Weizs\"acker--Williams field) & Eqs.~\eqref{eq:A1mom}, \eqref{eq:A1X}\\
$\mf A^{(3)}\equiv\mf A^{(3)}_2$ & emission of one gluon at order $g^3$; parts with one, two, three $\rho$ & Eqs.~\eqref{eq:A3onerho}--\eqref{eq:A3threerho}\\
$\mf B_2$ & emission of two gluons, order $g^2$ & Eqs.~\eqref{eq:B2mom}, \eqref{eq:B2X}\\
$\mf B_3$ & transition of a gluon into a gluon on the valence charges, order $g^2$ & Eqs.~\eqref{eq:B3mom}, \eqref{eq:B3X}\\
$\mf C\equiv\mf C_1$ & splitting of a gluon into two gluons, order $g$, independent of $\rho$ & Eqs.~\eqref{eq:C1mom}, \eqref{eq:C1X}\\
bar, e.g.\ $\bar{\mf A}$ & the same coefficient with every $\rho^a(\v u)\to U^{ab}(\v u)\rho^b(\v u)$ & Eq.~\eqref{eq:bar}\\[2pt]
\multicolumn{3}{@{}l}{\emph{Target}}\\
$S$ & eikonal $S$-matrix of the target & Eq.~\eqref{eq:Srotation}\\
$U^{ab}(\v x)$ & adjoint Wilson line; real, $U^\dagger=U^T$, $U^TU=1$ & Eq.~\eqref{eq:Srotation}\\
$\mathcal O_{\rm LO}$ & ${\rm Tr}\{[U^\dagger(\v x)-U^\dagger(\v w)][U(\v y)-U(\v w')]\}$ & Eq.~\eqref{eq:OLO}\\
$T^e$, $D^e$ & $(T^e)_{bc}=-if^{ebc}$, $(D^e)_{bc}=d^{ebc}$ & Eq.~\eqref{eq:TD}\\
$H_{\rm JIMWLK}$ & JIMWLK Hamiltonian; $d/dY=-H_{\rm JIMWLK}$ & Eq.~\eqref{eq:HJIMWLK}\\[2pt]
\multicolumn{3}{@{}l}{\emph{Transverse kernels}}\\
$\Kv^i(\v x)$ & $x^i/\v x^2$, Weizs\"acker--Williams field of a point charge & Eq.~\eqref{eq:kernels}\\
$K(\v x,\v y,\v z)$ & $\Kv(\v x-\v z)\cdot\Kv(\v y-\v z)$ & Eq.~\eqref{eq:kernels}\\
$K_D(\v x,\v y,\v z)$ & dipole kernel $K(\v x,\v x,\v z)-2K(\v x,\v y,\v z)+K(\v y,\v y,\v z)$ & Eq.~\eqref{eq:BFKLrhorho}\\
$\mathfrak K_{\rm row}$ & named kernel of a finite term of a row & Sec.~\ref{sec:fin_rows}\\
$\mathcal E(s)$, $\Psi(s)$ & recoil functions of Groups~I and~III & Eq.~\eqref{eq:recoilfns}\\[2pt]
\multicolumn{3}{@{}l}{\emph{Amplitudes and pieces of the result}}\\
$\mathcal A^a_i(k)$, $\mathcal A^{ab}_{ij}(k;q^+,\v v)$ & one- and two-gluon amplitudes & Eq.~\eqref{eq:amplitudes_def}\\
$\mathcal A_N,\dots,\mathcal A_{CB_2}$ & the five parts of the one-gluon amplitude of order $g^3$ & Eq.~\eqref{eq:A3parts}\\
$\mathcal A_B$, $\mathcal A^{(1,2)}_{AA}$, $\mathcal A_{CA}$ & the four parts of the two-gluon amplitude of order $g^2$ & Eq.~\eqref{eq:A2parts}\\
$\mathcal O(y,y')$, $\widehat L(\v q)$ & LO integrand; LO structure at transverse momentum $\v q$ & Eq.~\eqref{eq:dg_not}\\
$\mathcal N_0$, $\mathcal C$, $\mathcal C_{\rm f}$ & prefactors: $\frac{1}{(2\pi)^3}\frac{g^2}{2\pi^2}\frac1{k^+}$; $\frac{\alpha_s}{2\pi}\mathcal N_0$; $\mathcal C/\pi$ & Eqs.~\eqref{eq:dg_not}, \eqref{eq:Cf}\\
$\Pi(\xi)$, $\mathcal F(\xi)$ & LO structure of a parent of momentum $\v k/\xi$; the same with the collinear logarithm & Sec.~\ref{sec:dg_origin}, Eq.~\eqref{eq:dg_F}\\
$\mathcal X$ & BFKL structure of Row~II.IV & Eq.~\eqref{eq:dg_X}\\
$\mathcal P_{gg+}$, $\mathcal P_{gg}$, $\widetilde{\mathcal P}_{gg}$ & real splitting function; with $\frac{\beta_0}{2}\delta(1-\xi)$; with the additional subtraction at $\xi\to0$ & Eqs.~\eqref{eq:Pggplus}, \eqref{eq:DGLAP}\\
$\beta_0$, $\mathfrak a$ & $\frac{11}{3}N_c$ (pure Yang--Mills); $\alpha_sN_c/\pi$ & Eq.~\eqref{eq:dg_B}\\
\end{longtable}
\endgroup
The following terms are used throughout. A \emph{row} is one product of coefficients of $\bar\Omega^\dagger\Omega$ in the spectrum at order $g^4$ (Table~\ref{tab:rows}); rows are numbered within a Group, e.g.\ Row~I.III. The \emph{measured} gluon has momentum $k$ and stands at $\v w$ in the amplitude and at $\v w'$ in the conjugate amplitude; the second gluon, at $\v z$, is real and integrated over in Groups~II and~III and virtual in Group~I. Terms with $n$ charges are counted before any reordering; a commutator turns such a term into a term with $n-1$ charges at the same order in $g$. A finite row has up to three parts: the totally symmetrised three-$\rho$ product (``sym''), the two-$\rho$ contact terms left by the symmetrisation (``con''), and the genuine two-$\rho$ terms (``quad''). A \emph{tail} is a behaviour $\propto1/\v z^2$ of a term at large transverse distance of the integrated gluon, which alone would make the $\v z$ integral diverge logarithmically.
